\chapter{Introduction}
\label{ch:introduction}

\section{Motivation}
A useful answer often depends on a small part of a much larger document.
Finding that part is not just a matter of choosing the right document.
Even after the source is known, a retrieval system must decide what to
return from it. A phrase may identify a method, a paragraph may explain
its assumptions, and several passages may be needed to compare the
results of two experiments. The unit of retrieval therefore shapes what
an answering system can see.

Retrieval-augmented generation (RAG) brings retrieved material into the
generation process instead of relying only on information stored in a
model's parameters \cite{lewis2020rag}. This creates a useful separation
between acquiring evidence and producing an answer, but it also gives
the retrieval stage considerable responsibility. If relevant text is
missing, the generator has to work without it. If the supplied text is
mostly irrelevant, the generator must distinguish the useful evidence
from the surrounding material. Good retrieval is not sufficient for
good generation, but it is a necessary part of a dependable pipeline.

Chunking is one of the decisions that connects those two stages. In
fixed-size chunking, a document is divided into pieces according to a
token budget selected in advance. The same rule applies to all questions,
including those that need very different amounts of evidence. This is
convenient: indexing is straightforward, every chunk has a predictable
nominal size, and the retrieval procedure has few moving parts. It is
less clear whether that convenience comes at the expense of evidence
quality.

The intuition behind adaptive granularity is simple. Rather than ask
every question to work with the same chunk size, allow the system to
select a size that suits the current information need. The difficulty is
knowing what \enquote{suits} means before the answer is available. A
question may sound broad while its evidence is concentrated in one
sentence. A short factual question may depend on a table or on several
parts of a paper. Even a correct estimate of the total evidence length
does not specify how that evidence is distributed across retrievable
pieces.

Scientific papers make this problem particularly interesting. They
contain definitions, methodological explanations, experimental settings,
and results, often spread across separate sections. QASPER provides
questions anchored in such papers together with annotated supporting
evidence \cite{dasigi2021qasper}. This allows the granularity decision to
be studied against an explicit evidence target rather than against an
informal impression of what a passage should contain.

\section{Problem statement and scope}
This thesis asks whether an adaptive choice of chunk size can improve
evidence retrieval within a known source paper. The candidate sizes are
$\granularities=\{10,20,40,80,160\}$ GPT-2 tokens. The primary
question-level protocol selects one size and retrieves the five most
similar chunks at that size. Later experiments make the decision locally,
allowing different regions of the same paper to contribute chunks at
different sizes.

This is a deliberately controlled problem. The source paper is supplied
by the dataset, embeddings remain fixed, and retrieval uses cosine
similarity. The experiments do not train an embedding model, search an
unknown corpus for the correct paper, rerank passages with a cross-encoder,
or generate final answers. These choices isolate the granularity decision
and make its successes and failures easier to interpret.

The main outcome is joined evidence-token F1: the lexical overlap between
the combined retrieved text and the combined annotated evidence.
Chapter~\ref{ch:data} defines the metric precisely. Throughout the thesis,
\emph{retrieval F1} refers to this measure. It is distinct from the
classification F1 used to assess predicted size labels, and from the
answer-level scores used in a complete question-answering system.

The word \emph{adaptive} also needs a boundary. A method is adaptive here
when its size choice depends on the question or its similarity pattern
within the paper. An offline oracle that inspects gold evidence is not
a deployable adaptive method. It provides a reference for what could be
gained if the best of the available actions were known. Likewise, a
method that retrieves many more chunks is not automatically better
adaptation; it may simply spend a larger context budget.

\section{Research questions}
The investigation is organised around five research questions.

\begin{description}[style=nextline,leftmargin=0pt]
\item[RQ1: Is there useful headroom beyond a fixed granularity?]
How do the five fixed sizes compare, and how much can an offline
per-question choice improve over the strongest fixed setting?
\item[RQ2: Can the question alone support a useful size decision?]
Do question embeddings or a small language model predict a granularity
that improves evidence retrieval? How closely does label prediction
track the retrieval outcome?
\item[RQ3: Does the document's similarity structure add information?]
Can score distributions across coarse and fine spans help a router?
Does combining those signals with language-model features improve the
decision when stacking is trained without in-sample base predictions?
\item[RQ4: Is direct retrieval-utility supervision more effective?]
What changes when the target becomes the cost of a retrieval decision
rather than an evidence-length class?
\item[RQ5: Does local adaptation improve a matched retrieval pipeline?]
Can a system benefit from selecting different sizes in different regions,
and does a frozen language-model estimate of context need add useful
information to that local decision?
\end{description}

These questions reflect the actual development of the project. They are
not five independent confirmations of a single preselected method. The
answer to one question often motivated a change in the next experiment.
That history is important because it explains both the design choices
and the limits of the final evaluation.

\section{Experimental approach}
The work began with data ingestion and a fixed retrieval benchmark.
Documents were represented at all five sizes, embedded, and stored with
stable identifiers and source offsets. Questions and evidence were
normalised and audited. Fixed-size retrieval established the reference
outcomes from which oracle labels and later utility targets could be
constructed.

The first learned router used the question embedding. Logistic regression
and a small multilayer perceptron tested whether a comparatively cheap
representation was already sufficient. Mixed-size retrieval tested a
different possibility: perhaps it was unnecessary to predict a size at
all if candidates from all levels could share one ranking. These methods
provided practical baselines before introducing a language model.

The Qwen sequence then investigated the question-only decision in more
detail. It moved from frozen prompting to supervised numeric generation,
restricted alias outputs, class weighting, and a dedicated classification
head. A prompt change and a small learning-rate grid examined whether
the remaining limitation was primarily the output interface or the
training configuration.

The next sequence used the evidence landscape without using gold
evidence at inference. Similarity-tree features describe where high
question-to-chunk scores occur and whether they persist across scales.
Linear models, boosted trees, and Qwen-feature fusion were evaluated.
The fusion procedure was then corrected to use paper-grouped out-of-fold
base predictions, making the training interface closer to the way the
fusion model would encounter a new question.

Finally, expected-regret learning changed the objective, while local
routing changed the decision unit. A frozen short/medium/long Qwen feature
was added to the local router in the last experiment. Its unsuccessful
free-generation version and its constrained-output replacement are both
part of the record.

\section{Contributions}
The contribution is an evidence-based study of the granularity decision,
rather than a claim that one adaptive architecture universally solves it.
Four aspects of the work are particularly useful.

First, the repository provides a reproducible multi-granularity evaluation
setting with frozen questions, fixed embeddings, paper-restricted retrieval,
versioned metrics, and recoverable intermediate outputs. This makes it
possible to compare substantially different routers against the same
underlying retrieval actions.

Second, the study separates several meanings of an oracle. The
retrieval-optimal action, the total evidence-length category, and the
local evidence-overlap category supervise different decisions. Making
those differences explicit explains why classification improvements
sometimes move retrieval in the wrong direction.

Third, the experiments include methodological corrections and negative
results. In particular, the out-of-fold fusion rerun and the invalid-output
gate in the final Qwen feature experiment expose failure modes that would
be easy to hide in a report containing only the best scores.

Fourth, the expected-regret and matched local-routing comparisons identify
where adaptation became more competitive and where it still fell short.
The strongest learned question-level router reached 0.3075 mean retrieval
F1, but fixed 40-token retrieval reached 0.3159. Local routing similarly
did not demonstrate an advantage over its matched fixed-40 comparator.
These results support a more precise conclusion than either
\enquote{adaptation works} or \enquote{adaptation is unnecessary}.

\section{How to read the findings}
The repository's validation split was consulted across successive modelling
decisions. Although several later protocols use training-only selection,
prediction locks, and grouped folds, these safeguards cannot turn a
repeatedly observed validation split into an untouched test set. The
reported comparisons are therefore development findings. Their confidence
intervals quantify resampling uncertainty within the observed papers,
not all uncertainty due to model selection, random seeds, or transfer to
another domain.

This distinction does not make the experiments uninformative. It changes
what can be concluded from them. A carefully controlled improvement over
an earlier router can motivate the next study, while a final superiority
claim requires independent confirmation. The thesis keeps both ideas in
view: the results are useful, and their evidential scope is limited.

\section{Thesis organisation}
Chapter~\ref{ch:background} introduces the retrieval concepts and places
the work alongside research on granularity, hierarchy, and adaptive RAG.
Chapter~\ref{ch:data} describes QASPER, preprocessing, data integrity, and
evaluation targets. Chapter~\ref{ch:method} develops the baseline and
adaptive methods in experimental order. Chapter~\ref{ch:results} presents
the results, controlled comparisons, unsuccessful runs, and threats to
validity. Chapter~\ref{ch:conclusion} answers the research questions and
sets out the next experiments needed for a stronger conclusion. The
appendices document the chronology, configurations, output contracts,
and links between reported results and repository artifacts.
